\begin{lemma}
Let $1<p<\infty$, let $K\in\mathbb{R}$, let $\Omega\subseteq\mathbb{C}$ be open
and connected, and let $f\colon\mathbb{C}\to\mathbb{C}$ be $K$-quasiregular on
$\Omega$ in the sense of \noderef{QuasiregularOn}.  Let
$g\colon\mathbb{C}\to\mathbb{C}$ satisfy
\begin{enumerate}
\item[(i)] $g=f$ almost everywhere on $\Omega$;
\item[(ii)] $g$ is continuous on $\Omega$;
\item[(iii)] $g$ is real-differentiable at almost every point of $\Omega$;
\item[(iv)] $g$ solves equation (12) of the paper on $\Omega$ in the sense of
  \noderef{SolvesEq12}.
\end{enumerate}
Assume $g(z^{*})\neq0$ for some $z^{*}\in\Omega$, and let $z_0\in\Omega$
satisfy $g(z_0)=0$.  Assume the bundles (C2) and (C3) of \noderef{QCPackage}
and \noderef{SmoothOffZeros}.  Then there
are an integer $m\ge1$, a constant $K'\ge1$, an open set $V'$ with
$z_0\in V'\subseteq\Omega$ and maps $\chi,H\colon\mathbb{C}\to\mathbb{C}$ such
that
\begin{enumerate}
\item[(a)] $\chi$ is $K'$-quasiconformal on $V'$ in the sense of
  \noderef{QuasiconformalOn}, $\chi(z_0)=0$, and $g(z)=[\chi(z)]^{m}$ for every
  $z\in V'$;
\item[(b)] $\chi$ is $C^\infty$ on $V'-\{z_0\}$;
\item[(c)] $H(\chi(z))=z$ for every $z\in V'$, $H$ is continuous on
  $\chi(V')$, $H(0)=z_0$, and for every $\xi\in\chi(V')$ one has
  $H(\xi)\in V'$ and $\chi(H(\xi))=\xi$.
\end{enumerate}

This is the normalised form of the factorization (15) that Sections~3 to~5
use.  Clause (a) is the factorization theorem quoted in (C2), shrunk so that
$z_0$ is the only zero of $g$ in $V'$; (b) is \noderef{ChiSmoothOffZero}
applied on the punctured neighborhood, where (C3) makes $g$ smooth and the
discreteness of the zero set makes it nonvanishing; and (c) is the inverse
furnished by the quasiconformality of $\chi$, with the two inversion
identities and $H(0)=z_0$.

The factorization, the equation and the singular point are all read for the
continuous representative $g$ of $f$, not for $f$ itself, and they must be.
Quasiregularity as \noderef{QuasiregularOn} defines it constrains $f$ only
through $W^{1,2}_{\mathrm{loc}}$ membership and an almost-everywhere
inequality, both unchanged by altering $f$ on a null set; so no hypothesis on
$f$ alone can give a pointwise identity $f=\chi^{m}$ on an open set, and the
set $f^{-1}(0)$ is not determined modulo null sets, whereas $g^{-1}(0)$ is.
This is the same reading of the quoted (C2) and (C3) that \noderef{QCPackage}
and \noderef{SmoothOffZeros} carry --- the paper's own at paper.tex:156, where
the singular set is written $f^{-1}(0)=\{z\in\Omega:\nabla u(z)=0\}$ --- and a
single representative is threaded through all of them.

No smoothness of $H$ is claimed, and none is available here: the paper's arrow
runs $H\to\chi$ (paper.tex:416), and $H$'s regularity off the origin is
produced downstream, by term-by-term differentiation of the expansion (19) as
at paper.tex:387--388.  Claiming it here would invert the development, and it
is what forced the illegal fourth bundle in earlier cycles.  Nothing else about
$H$ is claimed either.  Its $W^{1,2}_{\mathrm{loc}}$ membership, almost
everywhere differentiability and $J_H>0$ almost everywhere are one application
of (C2)'s inverse clause away; and its quasiconformality --- which the paper
asserts flatly at paper.tex:179 but which (C2) does not quote, so that
paper.tex:113 requires it to be proved --- is assembled from those, from (18)
and from the homeomorphism clauses of $\chi$ in the node that applies
Corollary 1, which already carries the bundle (C2) and the $\chi$ data.
\end{lemma}
\begin{proof}
\emph{Step 1: the factorization and the isolated zero.}  The clause
\texttt{qr\_factorization} of \noderef{QCPackage}, applied with $K$, $\Omega$,
the pair $f,g$, the point $z_0$, the openness and connectedness of $\Omega$,
the quasiregularity of $f$, the hypotheses (i) and (ii) on $g$, the point
$z^{*}$ with $g(z^{*})\neq0$ and $g(z_0)=0$, provides an integer $m\ge1$, a
constant $K'$, an open set $V_0$ with $z_0\in V_0\subseteq\Omega$ and a map
$\chi$ with
\begin{equation}\tag{1}
  \chi\ \text{is }K'\text{-quasiconformal on }V_0,\qquad \chi(z_0)=0,\qquad
  g(z)=[\chi(z)]^{m}\ (z\in V_0).
\end{equation}
By \noderef{QuasiconformalOn} the first clause of (1) contains
$K'$-quasiregularity of $\chi$ on $V_0$, whose first clause is $1\le K'$ by
\noderef{QuasiregularOn}; so
\begin{equation}\tag{2}
  1\le m,\qquad 1\le K'.
\end{equation}
The clause \texttt{qr\_discrete\_zeros} of \noderef{QCPackage}, applied with
the same $K$, $\Omega$, the pair $f,g$, the openness and connectedness of
$\Omega$, the quasiregularity of $f$, the hypotheses (i) and (ii) on $g$ and
the point $z^{*}$, read at the point $z_0\in\Omega$, provides $r_0>0$ with
\begin{equation}\tag{3}
  g(z)\neq0\qquad\text{for every }z\in\Omega\ \text{with }z\neq z_0\
  \text{and}\ |z-z_0|<r_0 .
\end{equation}

\emph{Step 2: the neighborhood $V'$.}  Put
\[
  V'=V_0\cap\{z:|z-z_0|<r_0\} .
\]
It is open, being an intersection of two open sets; $z_0\in V'$, since
$z_0\in V_0$ and $|z_0-z_0|=0<r_0$; and $V'\subseteq V_0\subseteq\Omega$.  By
(3) and $V'\subseteq\Omega$,
\begin{equation}\tag{4}
  g(z)\neq0\qquad(z\in V'-\{z_0\}),
\end{equation}
and by (1), restricted to $V'\subseteq V_0$,
\begin{equation}\tag{5}
  \chi(z_0)=0,\qquad g(z)=[\chi(z)]^{m}\quad(z\in V').
\end{equation}

\emph{Step 3: $\chi$ is $K'$-quasiconformal on $V'$.}  By
\noderef{QuasiconformalOn} applied to $V_0$ we have: $\chi$ is
$K'$-quasiregular on $V_0$, injective on $V_0$, continuous on $V_0$, the image
$\chi(V_0)$ is open, and there is $\psi$ with $\psi(\chi(z))=z$ for $z\in V_0$
and $\psi$ continuous on $\chi(V_0)$.  We verify the same five clauses for
$V'$.

\emph{(3a) $K'$-quasiregularity on $V'$.}  By \noderef{QuasiregularOn} this
asks for $1\le K'$, which is (2), for $\chi\in W^{1,2}_{\mathrm{loc}}(V')$ in
the sense of \noderef{MemWLocOn}, and for a pair of weak Wirtinger derivatives
of $\chi$ on $V'$ satisfying the distortion inequality (11) of the paper almost
everywhere on $V'$.

For the last clause, let $a,b$ be weak Wirtinger derivatives of $\chi$ on
$V_0$, with underlying family $\chi_x,\chi_y$, such that
$|b|\le\kappa_0|a|$ almost everywhere on $V_0$, where
$\kappa_0=(K'-1)/(K'+1)$; these exist by \noderef{QuasiregularOn} applied on
$V_0$.  The same $a,b$ are weak Wirtinger derivatives of $\chi$ on $V'$: the
local integrability clause (a) of \noderef{IsWeakWirtingerOn} restricts,
because every compact subset of $V'$ is a compact subset of $V_0$ and
\noderef{MemLpLocOn} quantifies over those; the distributional-derivative
clause (b) restricts, because by \noderef{IsTestFn} the only clause of
``$\varphi$ is a test function for a set'' which mentions the set is
$\operatorname{tsupp}\varphi\subseteq V'\subseteq V_0$, so every test function
for $V'$ is one for $V_0$, while the integral identity of
\noderef{IsWeakPartialOn} does not mention the set; and clause (c) is a
pointwise identity, independent of the set.  The inequality also restricts,
since $V'\subseteq V_0$ and an almost-everywhere implication with a smaller
antecedent set is weaker.

For the Sobolev clause, take the family $D$
witnessing $\chi\in W^{1,2}_{\mathrm{loc}}(V_0)$: each $D(w)$ with $|w|\le1$
lies in $L^{2}_{\mathrm{loc}}(V')$ in the sense of \noderef{MemLpLocOn},
because every compact subset of $V'$ is a compact subset of $V_0$; and for
$|w|<1$ and $j\in\{0,1\}$ the identity of \noderef{IsWeakPartialOn} holds on
$V'$ because every test function for $V'$ is a test function for $V_0$ --- by
\noderef{IsTestFn} the only clause involving the set is
$\operatorname{tsupp}\varphi\subseteq V'\subseteq V_0$ --- and the identity
itself does not mention the set.  So $D$ witnesses
$\chi\in W^{1,2}_{\mathrm{loc}}(V')$.

\emph{(3b) injectivity and continuity on $V'$.}  Both restrict from $V_0$
along $V'\subseteq V_0$.

\emph{(3c) $\chi(V')$ is open.}  We claim
\begin{equation}\tag{6}
  \chi(V')=\chi(V_0)\cap\psi^{-1}(V') .
\end{equation}
If $\xi=\chi(z)$ with $z\in V'$, then $z\in V_0$, so $\xi\in\chi(V_0)$ and
$\psi(\xi)=\psi(\chi(z))=z\in V'$.  Conversely, let $\xi\in\chi(V_0)$ with
$\psi(\xi)\in V'$ and write $\xi=\chi(y)$ with $y\in V_0$; then
$\psi(\xi)=\psi(\chi(y))=y$, so $y\in V'$ and $\xi=\chi(y)\in\chi(V')$.  This
proves (6).  Now $\psi$ is continuous on $\chi(V_0)$ and $V'$ is open, so the
right-hand side of (6) is open in the subspace $\chi(V_0)$; and $\chi(V_0)$ is
open in $\mathbb{C}$.  A subset of $\mathbb{C}$ open in an open subspace is
open, so $\chi(V')$ is open.

\emph{(3d) the inverse clause.}  The same $\psi$ satisfies $\psi(\chi(z))=z$
for $z\in V'\subseteq V_0$ and is continuous on
$\chi(V')\subseteq\chi(V_0)$.

With (3a)--(3d) and (5), $\chi$ is $K'$-quasiconformal on $V'$ with
$\chi(z_0)=0$ and $g=\chi^{m}$ on $V'$, which is assertion (a).

\emph{Step 4: $g$ is $C^\infty$ on $V'-\{z_0\}$.}  By \noderef{QuasiregularOn}
the quasiregularity of $f$ on $\Omega$ contains
$f\in W^{1,2}_{\mathrm{loc}}(\Omega)$.  The clause
\texttt{smooth\_off\_zeros} of \noderef{SmoothOffZeros}, applied with $p$,
$\Omega$, the pair $f,g$, the openness of $\Omega$, $1<p$, that membership and
the four hypotheses (i)--(iv) on $g$, gives that $g$ is $C^\infty$ on
$\Omega-\{z:g(z)=0\}$.  By (4) every $z\in V'-\{z_0\}$ lies in $\Omega$ with
$g(z)\neq0$, so
\begin{equation}\tag{7}
  V'-\{z_0\}\subseteq\Omega-\{z:g(z)=0\},
\end{equation}
and a function of class $C^\infty$ on a set is of class $C^\infty$ on every
subset of it; hence $g$ is $C^\infty$ on $V'-\{z_0\}$.

\emph{Step 5: $\chi$ is $C^\infty$ on $V'-\{z_0\}$.}  The set $U=V'-\{z_0\}$
is open, $V'$ being open and $\{z_0\}$ closed.  The hypotheses of
\noderef{ChiSmoothOffZero} hold with this $m$, this $U$, and $g,\chi$ in the
roles of its two maps: $m\ge1$ is (2); (H1), continuity of $\chi$ on $U$, is
the restriction of (3b); (H2), $g=\chi^{m}$ on $U$, is (5) restricted to
$U\subseteq V'$; (H3), $g$ of class $C^\infty$ on $U$, is Step 4; and (H4),
$g\neq0$ on $U$, is (4).  Therefore
\begin{equation}\tag{8}
  \chi\ \text{is }C^\infty\ \text{on }U=V'-\{z_0\},
\end{equation}
which is assertion (b).

\emph{Step 6: the inverse $H$.}  Put $H=\psi$, the map of Step 3.  By (3d),
\begin{equation}\tag{9}
  H(\chi(z))=z\ \ (z\in V'),\qquad H\ \text{is continuous on }\chi(V') .
\end{equation}
Three consequences.  First, taking $z=z_0\in V'$ in (9) and using $\chi(z_0)=0$
from (5),
\begin{equation}\tag{10}
  H(0)=H\bigl(\chi(z_0)\bigr)=z_0 .
\end{equation}
Second, let $\xi\in\chi(V')$ and write $\xi=\chi(z)$ with $z\in V'$; then
$H(\xi)=z$ by (9), so
\begin{equation}\tag{11}
  H(\xi)\in V'\qquad\text{and}\qquad \chi\bigl(H(\xi)\bigr)=\chi(z)=\xi
  \qquad\bigl(\xi\in\chi(V')\bigr).
\end{equation}
With (9), (10) and (11), assertion (c) is proved.

\emph{Step 7: conclusion.}  The data $m$, $K'$, $V'$, $\chi$ and $H$ satisfy
$m\ge1$ and $K'\ge1$ by (2); $V'$ is open with $z_0\in V'\subseteq\Omega$ by
Step 2; assertion (a) holds by Step 3 with (5); (b) by (8); and (c) by (9),
(10) and (11).
\end{proof}
